\section{Which Paths Carry Evidence}
\label{sec:specificity}

\subsection{Only short paths are specific}
\label{sec:specificity-intact}

Table~\ref{tab:indications} and Figure~\ref{fig:summary}b show the intact full
graph. A shared protein links 31.0\% of the 6,233 approved indications and
14.2\% of their decoys, so a two-hop answer is 2.2 times more common among
indications. The contrast fades quickly with the horizon. At $L=3$ the ratio is
1.26, at $L=4$ it is 1.01, and without a limit both sets are connected in
99.97\% of cases. Near saturation the ratio of rates hides a residual difference: as odds
ratios, the separations are 2.71, 1.68 and 1.35 at $L=2$, 3 and 4. Mutual
information, which weighs both, falls from 29.6 to 11.9 and 0.7 millibits. Even
the two-hop signal is weak, 29.6 of a possible 1,000 millibits, or an AUROC of
0.59, which is what one should expect from a single shared protein. These ratios
are stable: over five further decoy draws the
two-hop ratio lies between 2.09 and 2.18, and a bootstrap over the 1,375
indicated drugs, which accounts for indications that share a drug, gives 95\%
intervals of 1.92--2.47 at $L=2$, 1.20--1.32 at $L=3$ and 1.00--1.02 at $L=4$.

\begin{table}[!htbp]
\centering
\small
\caption{Typed retrieval of 6,233 approved indications and 6,233
degree-preserving decoys on the intact full graph, with the chance rates of
Section~\ref{sec:theory-chance}: the configuration model with a Poisson
approximation (Eq.~\eqref{eq:cm}) and the bundled null model
(Eq.~\eqref{eq:ball}), evaluated for the decoy and the indication pairs.
Information is per pair, for a balanced mixture.}
\label{tab:indications}
\resizebox{\linewidth}{!}{%
\begin{tabular}{@{}lrrrrrrr@{}}
\toprule
& \multicolumn{3}{c}{Observed} & \multicolumn{3}{c}{Chance rate} & \\
\cmidrule(lr){2-4}\cmidrule(lr){5-7}
$L$ & Indications & Decoys & Ratio & CM, decoys & Null, decoys & Null, indications & Info (mbit) \\
\midrule
$L=2$      & 31.0\% & 14.2\% & 2.18 & 13.1\% & 13.4\% & 15.8\% & 29.6 \\
$L=3$      & 61.1\% & 48.3\% & 1.26 & 58.1\% & 49.1\% & 51.8\% & 11.9 \\
$L=4$      & 96.2\% & 94.9\% & 1.01 & 99.7\% & 95.3\% & 95.5\% & 0.7 \\
Unlimited  & 99.97\% & 99.97\% & 1.00 & -- & -- & -- & 0.0 \\
\bottomrule
\end{tabular}}
\end{table}

The chance models explain the decoys. The bundled null model,
Eq.~\eqref{eq:ball}, predicts the decoy rates at all three horizons to within
one percentage point, with no free parameter. Because decoys keep a real drug,
as the null model does, this confirms that a disease's protein links look random
with respect to the drug's neighborhood. The configuration model gets the
two-hop rate right (13.1\% against 14.2\%) and overestimates the three-hop rate
(58\% against 48\%). Its mean number of chance paths is roughly right (5.1
against 4.5 paths of length at most three per decoy), but the Poisson step is
not: chance paths come in bundles through the same drug target. Measured
against the same null model, indications exceed chance by a factor of 1.96 at
$L=2$, 1.18 at $L=3$ and 1.01 at $L=4$. The excess is an association present
in the curated data, and each extra hop buries it under chance paths, whose number
grows by a factor of order $\kappa_{\mathrm{PPI}}$ per step.

Equation~\eqref{eq:excess} condenses this into one number per horizon. The
observed rates give $r_L=2.42$, $1.43$ and $1.10$ at $L=2$, 3 and 4, so the
excess of real over chance paths, $r_L-1$, falls from 1.42 to 0.43 and 0.10: it
shrinks by a factor of 3.3 and then 4.5 (Figure~\ref{fig:frontier}b;
the four-hop value, $0.10\pm0.03$, is near the noise floor). The rate of this
decay can be read off the path counts. With one protein--protein step, the mean
number of chance paths per decoy grows from 0.26 to 4.25, a factor of 17, while
the excess of indications over decoys grows only from 0.50 to 1.97 paths, a
factor of 4. The specific signal does propagate, so indications carry excess
three-hop paths as network-proximity studies find \cite{guney2016network}, but
it propagates about four times more slowly than chance, consistent with
Section~\ref{sec:hubs}: real indications avoid the hub proteins through which
chance paths multiply. The ratio of the mean numbers of paths of length at most
$L$ then falls from 2.98 to 1.55, so its excess shrinks by a factor of 3.6,
close to the 3.3 of the effective excess. For the next hop the mean counts no longer predict it (a
factor of 1.9 against 4.5): almost every pair is connected at $L=4$, and the
counts are overdispersed (the variance of three-hop counts is 35--48 times
their mean), so path existence no longer tracks the mean. Up to $f=0.5$
damage changes $r_2$ and $r_3$ little (2.48 and 1.49), while the four-hop excess
halves ($r_4=1.05$); over this range the ratio of success rates rises mainly
because thinning lowers the chance rate. Beyond it $r_2$ itself grows (2.97 at
$f=0.8$) towards the ratio of mean path counts, 2.98, as heavy thinning makes
surviving paths rare, so that success follows the mean count rather than the
overdispersed distribution.

Could path counts at longer horizons recover what path existence loses? Among
pairs with no shared protein, the number of three-hop paths normalized by
$k_dk_y$ separates indications from decoys with an AUROC of only 0.53, and the
number of four-hop paths with 0.54, against 0.59 for the number of shared
proteins over all pairs. Our claims concern retrieval by path existence, the
rule that hop-limited retrieval uses, but counting longer paths does not change
the picture much. Contraindications sit between the two sets: a shared protein
links 20.2\% of the 21,835 contraindicated pairs, so a two-hop path signals a
mechanistic relation between a drug and a disease, not necessarily a
therapeutic one.

\subsection{The specific paths are the fragile ones}
\label{sec:specificity-damage}

The information sits in exactly the observable that Section~\ref{sec:theory}
identifies as fragile. Under random edge removal
(Figure~\ref{fig:summary}a), the giant component
still holds 86\% of nodes at $f=0.5$, and 85\% of all indications remain
connected without a hop limit. Two-hop recall falls from 31.0\% to 12.0\% (five orders). Of
the 1,931 two-hop indications, 40.3\% keep their link (mean of 500 orders), as
Eq.~\eqref{eq:r2} gives from their shared-protein counts (0.403). A single
damage state scatters around this mean with a standard deviation of 0.020,
twice the value for independent queries, because indications share drugs.

The separation between indications and decoys follows from the same formula.
Indications with a shared protein have more of them than decoys (2.45 against
1.80 on average; over all pairs, 0.76 against 0.26), so the ratio of the two success rates rises with damage, from
2.18 at $f=0$ towards $\langle m\rangle_{\mathrm{ind}}/\langle m\rangle_{\mathrm{dec}}=2.98$
as $f\to1$ (Section~\ref{sec:theory-r2}). Equation~\eqref{eq:r2} gives 2.44 at
$f=0.5$, against 2.39 measured. Fewer pairs have any answer, however, so the
information carried by a two-hop answer falls to 40\% of its intact value by
$f=0.5$ (Figure~\ref{fig:indications}). Applying Eq.~\eqref{eq:r2} to both sets
accounts for this decay to within 1 millibit at every damage level, inside the
spread of a single damage state. For $L=3$, the
independent-path estimate, Eq.~\eqref{eq:paths}, is a genuine
approximation and does almost as well (within 1.4 millibits); it underestimates the information at $f\ge0.7$, where a
three-hop answer carries about as much information as a two-hop one (3.9 against
3.5 millibits at $f=0.8$, within one standard deviation). The advantage of short
answers therefore holds over most of the damage range and disappears only when
the graph has lost most of its facts.

\begin{figure}[!htbp]
  \centering
  \includegraphics[width=0.66\linewidth]{act6/fig_information.pdf}
  \caption{Mutual information between the label (indication or decoy) and the
  retrieval outcome on the full graph under random edge removal, in millibits
  per pair, for each hop limit. Shading shows one standard deviation over five
  damage orders. Black lines follow from Eq.~\eqref{eq:r2} for $L=2$ (exact in
  expectation) and from the independent-path estimate, Eq.~\eqref{eq:paths}, for
  $L=3$.}
  \label{fig:indications}
\end{figure}
